\section{Time-Dependent Linear Response Theory in Noisy Network Dynamics}\label{sec:time-dependent}
\quad The steady-state response developed above characterizes the difference between two nearby stationary states before and after a small parameter perturbation. In many situations, however, a perturbation is applied through a finite-time protocol rather than instantaneously. We therefore consider parameter protocols that drive the system from one steady state to another. Starting from a system that initially fluctuates near a fixed point, we derive a time-dependent response theory that predicts the average change of observables as a function of time. In particular, we derive the responses of the mean state, deterministic force, and equal-time covariance to time-dependent perturbations of the network parameters. We then validate the theory numerically for linear network dynamics in both an equilibrium reference network and a nonequilibrium steady-state network.
\subsection{General Formula for Time-Dependent Linear Response Function}
\quad We consider the response induced by a time-dependent perturbation of the network parameters, $\vec{p}\rightarrow\vec{p}+\delta \vec{p}(t)$, and assume that the system is in the unperturbed steady state for $t\le 0$. The time-dependent linear-response function of an observable $\mathbf{B}(t)$ is defined as
\begin{equation}
    \langle \mathbf{B}(t) \rangle_{\vec{p}+\delta \vec{p}(t)} = \langle \mathbf{B}(0) \rangle_{ss,\vec{p}} + \int_0^t dt' \boldsymbol{\chi}(t-t') \delta \dot{\vec{p}}(t').
    \label{time-dependent_response}
\end{equation}
The conjugate force $\vec{\mathscr{F}}(t)$ is defined in the same way as in Eq. (\ref{conjugate_force}). The time-dependent linear-response function is then given by
\begin{eqnarray}
    \boldsymbol{\chi}(t) &=&\langle \nabla_p{\mathbf{B}} \rangle_{ss,{\vec p}} + \int_0^t dt'\langle {\mathbf {B}}(t-t'){\dot{\vv {\mathscr F}}}^\intercal\rangle_{ss,{\vec p}} \nonumber\\
    &=& \boldsymbol{\chi} ^{(ss)}- \langle \delta{\mathbf B}(t)\delta{\vv {\mathscr F}}^\intercal(0)\rangle_{ss,{\vec p}}=\boldsymbol{\chi} ^{(ss)}- \langle \delta{\mathbf B}(t){\vv {\mathscr F}}^\intercal(0)\rangle_{ss,{\vec p}}.
    \label{chineqm}
\end{eqnarray}
The corresponding time-dependent response of $\mathbf{B}(t)$ is
\begin{equation}
    \langle \mathbf{B}(t)\rangle_{{\vec p}+\delta{\vec p}(t)}= \langle \mathbf{ B}(0)\rangle_{ss,{\vec p}}+\boldsymbol{\chi}^{(ss)}\delta {\vec p}(t)-\int_0^t dt'\langle \delta\mathbf{ B}(t-t'){\vv {\mathscr F}}^\intercal(0)\rangle_{ss,{\vec p}} \delta{\dot {\vec p}}(t').
    \label{Bnoneqm}
\end{equation}
For the linearized noisy dynamics, both the time-lagged covariance $\mathbf{K}_\tau$ and the mixed correlation with the conjugate force are governed by $e^{\tau \mathbf{Q}}$. Using the time-dependent conjugate-force relation derived in Appendix~B, one obtains
\begin{equation}
    \langle \delta x_i(\tau)\mathscr{F}_\alpha(0)\rangle = \sum_\ell \frac{\partial X_\ell}{\partial p_\alpha}
    (e^{\tau \mathbf{Q}})_{i\ell}.
    \label{time_mixed_xF}
\end{equation}

This result shows that, within the linearized theory, the time-dependent response is determined by two ingredients: the parameter-induced shift of the noise-free fixed point and the propagator $e^{\tau\mathbf{Q}}$, which governs how the perturbation is transmitted through the network in time.
\subsection{Response Functions in the Linearized-Dynamics Framework}
\quad We now specialize the general time-dependent response formulas to the linearized dynamics around the stable fixed point. The response functions below follow from Eqs. (\ref{chineqm}) and (\ref{time_mixed_xF}) within this linearized Gaussian framework. Eqs. (\ref{chinsslinear})--(\ref{chitKolinear}) assume the local linearization around the fixed point.\par
For the mean state, combining Eq. (\ref{chineqm}) with Eq. (\ref{time_mixed_xF}) gives
\begin{equation}
    \chi_{i\alpha}^{(x)}(t)=\left(e^{t{\bf Q}}{\bf Q}^{-1} \nabla_p{\vec { F}}|_{\vec X} \right)_{i\alpha}-\left({\bf Q}^{-1} \nabla_p{\vec { F}} |_{\vec X}\right)_{i\alpha}. \label{chinsslinear}
\end{equation}
For the deterministic force, the linearized fluctuation satisfies $\delta \vec{F}(t)=\mathbf{Q}\,\delta \vec{x}(t)$ to leading order. The corresponding time-dependent response function is therefore
\begin{equation}
    \chi_{i\alpha}^{(F)}(t)=\left(e^{t{\bf Q}}\nabla_p{\vec { F}}|_{\vec X} \right)_{i\alpha}.
    \label{chitforcelinear}
\end{equation}
For the equal-time covariance matrix $\mathbf{K}(t)\equiv \langle \delta \vec{x}(t) \delta \vec{x}^\intercal(t) \rangle$, its time-dependent response function gives
\begin{equation}
    \chi_{ij\alpha}^{(K)}(t) = \chi_{ij\alpha}^{(ss)} - \sum_{\ell,m} \left((e^{t\mathbf{Q}})_{i\ell}\chi_{\ell m\alpha}^{(ss)}(e^{t\mathbf{Q}^\intercal})_{mj}\right),
    \label{chitKolinear}
\end{equation}
where $\chi_{\ell m\alpha}^{(ss)}$ is the steady-state response function derived in Section IV. These expressions provide the explicit time-dependent response kernels within the linearized framework. We next specify the simulation setting and then evaluate them further for the particular linear force and diffusive coupling used in the numerical tests.
\subsection{Simulation Setup and Validation Strategy}
To validate the explicit time-dependent formulas derived for linearized network dynamics, we consider the two four-node networks shown in Fig. \ref{graphN4}: an undirected network satisfying the equilibrium condition and a directed network realizing a nonequilibrium steady state. In both cases, the node parameters and nonzero coupling weights are fixed at $r_i=5$ and $W_{ij}=5$, respectively. The intrinsic dynamics is chosen as
\begin{equation}
    f_i(x_i)=r_i(a_i-x_i),
\end{equation}
with intrinsic target states $\vec{a}=(2,4,6,8)$, and the coupling function is diffusive,
\begin{equation}
    h(x_i,x_j)=x_j-x_i .
\end{equation}
For the undirected network, all nodes have identical noise strengths, $\sigma_{ii}=0.01$, so the system serves as an equilibrium reference case. In contrast, the directed network contains four unidirectional edges, thereby producing a nonequilibrium steady state.\par
We perturb either the node parameter $r_1(t)$ or a single coupling weight $W_{12}(t)$, corresponding to the edge from node 2 to node 1 under our convention. For the undirected network, the same perturbation is applied to the reciprocal edge between nodes 1 and 2. The time-dependent perturbation protocols are given in Eqs. \eqref{dra(t)cos} and \eqref{dWab(t)cos}, with amplitude $b=0.5$ and protocol duration $\tau=10$. The same protocol parameters are used for all time-dependent comparisons shown below. In the following subsections, we evaluate the corresponding explicit linear-response formulas for single-node and single-edge perturbations and compare them with direct Langevin simulations.

The observables used for validation are the mean state $\langle x_i(t)\rangle$, 
the mean force $\langle F_i(t)\rangle$, and the equal-time covariance matrix
\begin{equation}
    K_{ij}(t,t) = \left\langle \delta x_i(t)\delta x_j(t)\right\rangle .
\end{equation}
All ensemble averages are computed from $10^5$ independent trajectories over the time interval $0\leq t\leq 20$. Before the perturbation is applied at $t=0$, the system is first allowed to relax to its unperturbed steady state. In the time-series comparisons below, dashed curves denote direct Langevin simulations and solid curves denote the linear-response prediction. Since showing all components would be visually redundant for this four-node example, we display the representative quantities $\langle x_1(t)\rangle$, $\langle x_2(t)\rangle$, $\langle F_1(t)\rangle$, $\langle F_2(t)\rangle$, $K_{11}(t,t)$, and $K_{12}(t,t)$. The remaining components show similarly good agreement, so only representative responses are plotted.

\begin{figure}[H]
    \centering
    \subfigure[]{\includegraphics*[width=.45\columnwidth]{Figure/Time_dependent/undirected_graph_N4.eps}}
    \subfigure[]{\includegraphics*[width=.45\columnwidth]{Figure/Time_dependent/directed_graph_N4.eps}} 
    \caption{Four-node networks used to validate the time-dependent response formulas. (a) Undirected cycle with uniform noise, which serves as an equilibrium reference case. (b) Directed cycle with heterogeneous noise, which realizes a nonequilibrium steady state. In the numerical tests, we perturb either node $1$ through $\delta r_1(t)$ or the directed edge from node $2$ to node $1$ through $\delta W_{12}(t)$; for the undirected network, the reciprocal edge is changed by the same amount to preserve the symmetric topology.}
    \label{graphN4} 
\end{figure}

\subsection{Time-dependent Response due to the time-dependent perturbation in a node parameter}
\quad We first consider a time-dependent perturbation of a single node parameter $\delta r_\alpha(t)$. For $f_i(x_i;r_i) = r_i(a_i-x_i)$ and $h(x_i,x_j)=x_j-x_i$, the time-dependent response of $x_i$ is
\begin{equation}
    \langle x_i(t)\rangle_{r+\delta r_\alpha(t)}=X_i+ (a_\alpha-X_\alpha)\bigg[\int_0^tdt'\left(e^{(t-t'){\bf Q}}{\bf Q}^{-1}  \right)_{i\alpha}\delta{\dot r}_\alpha(t') -Q^{-1}_{i\alpha} \delta r_\alpha(t)\bigg].\label{xi(t)drlinear}
\end{equation}
In the numerical tests, we use the cosine protocol
\begin{equation}
     \delta r_\alpha(t)=\begin{cases} \frac{b}{2}(1-\cos\frac{\pi t}{\tau}), & 0\leq t\leq\tau\\
     b &t>\tau
     \end{cases},\label{dra(t)cos}
\end{equation}
As $t\to\infty$, since the eigenvalues of $\mathbf{Q}$ are negative, $\langle x_i(t)\rangle$ approaches the new noise-free fixed point. Using Eq. (\ref{chitforcelinear}), the mean-force response is
\begin{equation}
    \langle F_i(t)\rangle_{r+\delta r_\alpha(t)}
    =(a_\alpha-X_\alpha)\int_0^t dt'
    \left(e^{(t-t'){\bf Q}}\right)_{i\alpha}\delta{\dot r}_\alpha(t').
    \label{Fi(t)drlinear}
\end{equation}
The ensemble averages of $\vec{x}(t)$ and $\vec{F}(t)$ do not depend on the noise strengths and are determined solely by ${\bf Q}$. The equal-time covariance response is
\begin{equation}
    K_{ij}(t,t)=({\bf K_0})_{ij}
    -\left( \frac{\partial {\bf K}_0}{\partial Q_{\alpha\alpha}} \right)_{ij}\delta r_\alpha(t)
    +\int_0^t dt' \left( e^{(t-t'){\bf Q}}\frac{\partial {\bf K}_0}{\partial Q_{\alpha\alpha}} e^{(t-t'){\bf Q}^\intercal} \right)_{ij}\delta \dot{r}_\alpha(t').
    \label{Kij(t)drlinear}
\end{equation}
where $\mathbf{K}_0$ denotes the steady-state covariance for $t\leq 0$. For the cosine protocol (\ref{dra(t)cos}), the convolution integrals in Eqs. (\ref{xi(t)drlinear})--(\ref{Kij(t)drlinear}) can be analyzed by spectral decomposition of ${\bf Q}$. Appendix C gives the resulting closed-form expressions for $\langle x_i(t)\rangle$ and $\langle F_i(t)\rangle$, together with the mode-expanded covariance response written in terms of the protocol integral $I_{m\ell}(t)$.
\begin{figure}[H]
    \centering
    \subfigure[]{\includegraphics*[width=.45\columnwidth]{Figure/Time_dependent/nonsteady_response_xunddr.eps}}
    \subfigure[]{\includegraphics*[width=.47\columnwidth]{Figure/Time_dependent/nonsteady_response_funddr.eps}} 
    \subfigure[]{\includegraphics*[width=.45\columnwidth]{Figure/Time_dependent/nonsteady_response_Kounddr.eps}}
    \caption{Comparison between Langevin simulations (dashed) and the time-dependent linear-response prediction (solid) for a node perturbation $\delta r_1(t)$ in the four-node undirected network of Fig. \ref{graphN4}a. (a) Mean-state response $\langle x_i(t)\rangle-\langle x_i(0)\rangle$ for the components $i=1,2$. (b) Mean-force response $\langle F_i(t)\rangle$ for $i=1,2$. (c) Equal-time covariance response for the components $K_{11}(t,t)$ and $K_{12}(t,t)$. The agreement over both the driving stage and the post-driving relaxation verifies the time-dependent response theory in the equilibrium case.}
    \label{Time-dependent_undirected_dr} 
\end{figure}
Figures \ref{Time-dependent_undirected_dr} and \ref{Time-dependent_directed_dr} compare these predictions with Langevin simulations for a perturbation of the node parameter $r_1(t)$. The curves resolve both stages of the response: the driven interval $0\le t\le \tau$, during which the observables follow the smooth protocol, and the subsequent relaxation for $t>\tau$, during which the system approaches the new steady state. Agreement for $\langle x_i(t)\rangle$ and $\langle F_i(t)\rangle$ verifies the predicted mean relaxation dynamics. Agreement for $K_{ij}(t,t)$ shows that the theory also captures the time evolution of fluctuations and correlations.
\begin{figure}[H]
    \centering
    \subfigure[]{\includegraphics*[width=.45\columnwidth]{Figure/Time_dependent/nonsteady_response_xdirdr.eps}}
    \subfigure[]{\includegraphics*[width=.47\columnwidth]{Figure/Time_dependent/nonsteady_response_fdirdr.eps}} 
    \subfigure[]{\includegraphics*[width=.45\columnwidth]{Figure/Time_dependent/nonsteady_response_Kodirdr.eps}}
    \caption{Same comparison as in Fig. \ref{Time-dependent_undirected_dr}, but for the directed network in Fig. \ref{graphN4}b, which is in a nonequilibrium steady state. The observables are again (a) $\langle x_1(t)\rangle-\langle x_1(0)\rangle$ and $\langle x_2(t)\rangle-\langle x_2(0)\rangle$, (b) $\langle F_1(t)\rangle$ and $\langle F_2(t)\rangle$, and (c) $K_{11}(t,t)$ and $K_{12}(t,t)$. The close overlap between theory and simulation shows that the time-dependent response formulas remain valid even when detailed balance is broken.}
    \label{Time-dependent_directed_dr} 
\end{figure}
Increasing $r_1$ strengthens the local restoring force on node 1 toward its intrinsic target $a_1$. Since the unperturbed state satisfies $X_1>a_1$ because of the couplings, the perturbation drives $\langle x_1(t)\rangle$ downward and gives a negative response in $\langle F_1(t)\rangle$. The stronger local relaxation also suppresses fluctuations around node 1, leading to a decrease in $K_{11}(t,t)$. In the undirected case, the perturbation is transmitted more directly to node 2, so both the coupled response and the covariance $K_{12}(t,t)$ decrease. In the directed case, node 2 is not directly affected by node 1, and its response can therefore differ in amplitude and transient shape.
\subsection{Time-dependent Response due to the time-dependent perturbation in a connection weight}
\quad We next consider a time-dependent perturbation of a single connection weight. The derivation parallels the node-perturbation case, so we state the corresponding formulas for the same linear force and diffusive coupling used in the simulations and then compare them with direct simulations.
For this specific network model, the relaxation response of the mean position is
\begin{equation}
    \langle x_i(t)\rangle_{{\bf W}+\delta W_{\alpha\beta}(t)}
    =X_i+ (X_\beta-X_\alpha)\bigg[\int_0^t dt'\left(e^{(t-t'){\bf Q}}{\bf Q}^{-1} \right)_{i\alpha}\delta{\dot W}_{\alpha\beta}(t') -Q^{-1}_{i\alpha} \delta W_{\alpha\beta}(t)\bigg].
    \label{xi(t)dWlinear}
\end{equation}
We again use the cosine protocol on the connection between node $\alpha$ and node $\beta$:
\begin{equation}
     \delta W_{\alpha\beta}(t)=\begin{cases} \frac{b}{2}(1-\cos\frac{\pi t}{\tau}), & 0\leq t\leq\tau\\
     b &t>\tau
     \end{cases},\label{dWab(t)cos}
\end{equation}
The mean-force response is
\begin{equation}
    \langle F_i(t)\rangle_{{\bf W}+\delta W_{\alpha\beta}(t)}
    =(X_\beta-X_\alpha)\int_0^t dt'
    \left(e^{(t-t'){\bf Q}}\right)_{i\alpha}\delta{\dot W}_{\alpha\beta}(t').
    \label{Fi(t)dWlinear}
\end{equation}
The equal-time covariance response is
\begin{eqnarray}
    K_{ij}(t,t)&=&({\bf K_0})_{ij}
    +\left[- \frac{\partial {\bf K}_0}{\partial Q_{\alpha\alpha}}+\frac{\partial {\bf K}_0}{\partial Q_{\alpha\beta}} \right]_{ij}\delta W_{\alpha\beta}(t)\nonumber\\
    &+&\int_0^t dt' \left[ e^{(t-t'){\bf Q}}\left(-\frac{\partial {\bf K}_0}{\partial Q_{\alpha\alpha}} +\frac{\partial {\bf K}_0}{\partial Q_{\alpha\beta}}\right) e^{(t-t'){\bf Q}^\intercal} \right]_{ij} \delta \dot{W}_{\alpha\beta}(t').
    \label{Kij(t)dWlinear}
\end{eqnarray}
For the cosine protocol (\ref{dWab(t)cos}), the convolution integrals in Eqs. (\ref{xi(t)dWlinear})--(\ref{Kij(t)dWlinear}) can be evaluated explicitly by the same spectral decomposition of ${\bf Q}$ used in Appendix C. Figures \ref{Time-dependent_undirected_dW} and \ref{Time-dependent_directed_dW} show the corresponding validation for a time-dependent edge perturbation. The response is driven by the contrast $X_\beta-X_\alpha$ across the perturbed link, so the transient signals directly probe how a local change in coupling propagates through the network and modifies both the deterministic relaxation and the covariance dynamics.\par
Increasing $W_{12}$ strengthens the pull of node 2 on node 1. Since the unperturbed state satisfies $X_2>X_1$, the perturbation drives $\langle x_1(t)\rangle$ upward and gives a positive response in $\langle F_1(t)\rangle$. In the undirected case, the same perturbation also acts on $W_{21}$, reduces the state difference between the two nodes, and enhances their correlation, so $K_{12}(t,t)$ increases while $K_{11}(t,t)$ decreases. In the directed case, the response of node 2 is not directly affected by $\delta W_{12}$, so its amplitude is smaller and its time dependence is set by the network propagation.
\begin{figure}[H]
    \centering
    \subfigure[]{\includegraphics*[width=.45\columnwidth]{Figure/Time_dependent/nonsteady_response_xunddW.eps}}
    \subfigure[]{\includegraphics*[width=.47\columnwidth]{Figure/Time_dependent/nonsteady_response_funddW.eps}} 
    \subfigure[]{\includegraphics*[width=.45\columnwidth]{Figure/Time_dependent/nonsteady_response_KounddW.eps}}
    \caption{Comparison between Langevin simulations (dashed) and the time-dependent linear-response prediction (solid) for an edge perturbation in the four-node undirected network. The coupling on the link between nodes $1$ and $2$ is varied according to $\delta W_{12}(t)=\delta W_{21}(t)$. (a) Mean-state response for $\langle x_1(t)\rangle-\langle x_1(0)\rangle$ and $\langle x_2(t)\rangle-\langle x_2(0)\rangle$. (b) Mean-force response for $\langle F_1(t)\rangle$ and $\langle F_2(t)\rangle$. (c) Equal-time covariance response for $K_{11}(t,t)$ and $K_{12}(t,t)$. The agreement confirms the predicted transient and long-time response for a time-dependent perturbation of an undirected coupling.}
    \label{Time-dependent_undirected_dW} 
\end{figure}
\begin{figure}[H]
    \centering
    \subfigure[]{\includegraphics*[width=.45\columnwidth]{Figure/Time_dependent/nonsteady_response_xdirdW.eps}}
    \subfigure[]{\includegraphics*[width=.47\columnwidth]{Figure/Time_dependent/nonsteady_response_fdirdW.eps}} 
    \subfigure[]{\includegraphics*[width=.45\columnwidth]{Figure/Time_dependent/nonsteady_response_KodirdW.eps}}
    \caption{Same comparison as in Fig. \ref{Time-dependent_undirected_dW}, but for a directed perturbation $\delta W_{12}(t)$ in the nonequilibrium network of Fig. \ref{graphN4}b. Panels (a)-(c) show components of the mean state, mean force, and equal-time covariance, respectively. The close agreement between simulation and prediction demonstrates that the time-dependent response theory also captures the propagation of directed coupling perturbations in a nonequilibrium steady state.}
    \label{Time-dependent_directed_dW} 
\end{figure}
\subsection{Time-dependent Response due to the time-dependent perturbation in a noise strength}
\quad We next consider a time-dependent perturbation of the diagonal noise strength, $\delta \sigma_{\alpha\alpha}(t)$. Unlike node and edge perturbations, a perturbation in $\sigma_{\alpha\alpha}$ does not modify the deterministic drift matrix ${\bf Q}$. Therefore, for linear dynamics with additive zero-mean noise, the ensemble averages of the mean state and deterministic force remain unchanged to first order. In other words, a perturbation of the noise strength affects only the fluctuation sector at linear order, while the nontrivial response appears in the equal-time covariance matrix. Using Eq.~(\ref{Bnoneqm}), the time-dependent covariance response is
\begin{equation}
    {\bf K}(t,t)
    =
    {\bf K}_0
    +
    \frac{\partial {\bf K}_0}{\partial \sigma_{\alpha\alpha}}
    \,\delta \sigma_{\alpha\alpha}(t)
    -
    \int_0^t dt'\,
    e^{(t-t'){\bf Q}}
    \frac{\partial {\bf K}_0}{\partial \sigma_{\alpha\alpha}}
    e^{(t-t'){\bf Q}^\intercal}
    \,\delta \dot{\sigma}_{\alpha\alpha}(t').
    \label{Ksigma_time_response}
\end{equation}
This expression has the same structure as in the node- and edge-perturbation cases, except that the kernel is now determined by $\partial {\bf K}_0/\partial \sigma_{\alpha\alpha}$ rather than by derivatives with respect to entries of ${\bf Q}$. With the spectral decomposition ${\bf Q}=\sum_m \lambda_m {\vec u_m}{\vec v_m}^\intercal$, the covariance response can be written as
\begin{equation}
    K_{ij}(t,t)
    =
    (K_0)_{ij}
    +
    \left(
    \frac{\partial {\bf K}_0}{\partial \sigma_{\alpha\alpha}}
    \right)_{ij}
    \delta \sigma_{\alpha\alpha}(t)
    -
    \sum_{m,\ell}
    u_{im}u_{j\ell}
    \left(
    {\vec v_m}^\intercal
    \frac{\partial {\bf K}_0}{\partial \sigma_{\alpha\alpha}}
    {\vec v_\ell}
    \right)
    I_{m\ell}(t),
    \label{Ksigma_mode_expansion}
\end{equation}
where the protocol integral $I_{m\ell}(t)$ is the same as that defined in Eq.~(\ref{Iml}), with $\delta r_\alpha(t)$ replaced by $\delta \sigma_{\alpha\alpha}(t)$. For the cosine protocol used in the numerical test, we apply a time-dependent perturbation to $\sigma_{11}$ as
\begin{equation}
    \delta \sigma_{\alpha\alpha}(t)=\begin{cases} \frac{b}{2}(1-\cos\frac{\pi t}{\tau}), & 0\leq t\leq\tau \\
              b                                      & t>\tau
    \end{cases}.\label{dsaa(t)cos}
\end{equation}
Figure \ref{Time-dependent_Ko_ds} compares the predicted and simulated time-dependent responses of $K_{11}(t,t)$ and $K_{12}(t,t)$ for the two four-node graphs in Fig. \ref{graphN4}. Since the perturbation acts directly on the noise strength at node $1$, the most pronounced effect appears in the local variance $K_{11}(t,t)$. The cross-covariance $K_{12}(t,t)$ responds more indirectly through the propagation of fluctuations across the network. In the directed network of Fig. \ref{graphN4}b, node $2$ is not directly driven by node $1$, so the response of $K_{12}(t,t)$ is weaker than in the undirected case. In contrast, the bidirectional couplings in Fig. \ref{graphN4}a transmit the noise perturbation more efficiently, leading to a stronger covariance response.
\begin{figure}[H]
    \centering
    \subfigure[]{\includegraphics*[width=.45\columnwidth]{Figure/Time_dependent/nonsteady_response_Ko11ds.eps}}
    \subfigure[]{\includegraphics*[width=.45\columnwidth]{Figure/Time_dependent/nonsteady_response_Ko12ds.eps}}
    \caption{Comparison between Langevin simulations (dashed) and the time-dependent linear-response prediction (solid) for a noise-strength perturbation $\delta \sigma_{11}(t)$ in the two four-node networks of Fig. \ref{graphN4}. Panel (a) shows the equal-time variance $K_{11}(t,t)$ and panel (b) shows the equal-time covariance $K_{12}(t,t)$. In each panel, the responses for the undirected equilibrium network [Fig. \ref{graphN4}a] and the directed nonequilibrium network [Fig. \ref{graphN4}b] are plotted together. The close agreement confirms that a time-dependent perturbation of the local noise strength is captured accurately by the covariance response formula, while the weaker $K_{12}(t,t)$ signal in the directed case reflects the reduced propagation of fluctuations from node $1$ to node $2$.}
    \label{Time-dependent_Ko_ds}
\end{figure}
